\begin{frame}{Ontología de campos relativistas}
  \begin{itemize}
    \item \textbf{Campo espinorial continuo.} Excitación en el álgebra del
    espaciotiempo ($\mathcal{C}\ell_{1,3}$).
    \item \textbf{Evolución covariante.} Regida por el operador de Dirac
    $\nabla = \gamma^\mu \nabla_\mu$.
    \item \textbf{Régimen local del microscopio.} Curvatura gravitacional
    nula; geometría plana de Minkowski ($\eta_{\mu\nu}$).
    \item \textbf{Límite de interacción.} Reducido al acoplamiento mínimo
    con el potencial vector electromagnético ($A_\mu$).
  \end{itemize}
\end{frame}

\begin{frame}{Límite no relativista y atenuación del vacío}
  \begin{itemize}
    \item \textbf{Recuperación efectiva del campo.} Factorizar las fases de
    masa en reposo produce el campo de Schrödinger-Pauli.
    \item \textbf{Anomalía de potencial.} La brecha actúa como barrera
    clásicamente prohibida ($V > E$).
    \item \textbf{Valores propios imaginarios.} El operador momentum proyecta
    valores puramente imaginarios en la brecha.
    \item \textbf{Penetración del campo.} Cese de oscilaciones de onda plana;
    atenuación exponencial estricta en la amplitud.
  \end{itemize}
\end{frame}

\begin{frame}{Corrientes de Noether y flujo de probabilidad}
  \begin{itemize}
    \item \textbf{Superposición asintótica.} La proximidad subnanométrica
    fuerza la superposición de colas con decaimiento exponencial.
    \item \textbf{Transporte conservado.} Transporte de la corriente de
    probabilidad de Noether $J^\mu = \Psi \gamma^\mu \tilde{\Psi}$.
    \item \textbf{Ruptura de simetría.} El sesgo macroscópico ($V$)
    desalinea los niveles de Fermi.
    \item \textbf{Hamiltoniano de transferencia.} Media un flujo continuo
    desde estados ocupados hacia estados vacíos.
  \end{itemize}
\end{frame}

\begin{frame}{Observables y colapso de Tersoff-Hamann}
  \begin{itemize}
    \item \textbf{Paradigma de medición.} Evaluación de convoluciones
    espaciales del campo, no topografía física.
    \item \textbf{Modelo de Tersoff-Hamann.} El límite de la sonda se modela
    como un orbital $s$ perfectamente simétrico.
    \item \textbf{Colapso integral.} La simetría esférica reduce las
    superposiciones a evaluaciones puntuales en el ápice.
    \item \textbf{Proporcionalidad central.} La corriente mapea la densidad
    local de estados $I \propto \rho_S(E_F, \vec{r}_0)$.
  \end{itemize}
\end{frame}

\begin{frame}{Espacio de Fock y anisotropía geométrica}
  \begin{itemize}
    \item \textbf{Espacio de Fock.} Redes macroscópicas requieren
    distribuciones valoradas en operadores.
    \item \textbf{Sondas anisotrópicas.} Puntas funcionalizadas trascienden
    el acoplamiento de amplitud escalar.
    \item \textbf{Acoplamiento diferencial.} Matriz de transición acoplada a
    derivadas geométricas del campo ($\nabla \Psi$).
    \item \textbf{Resolución mejorada.} Detección de curvatura abstracta
    invisible para sondas de simetría esférica.
  \end{itemize}
\end{frame}
